\documentclass[letterpaper,12pt]{article}
%\usepackage{harvard}
%\usepackage{subfigure}
%\usepackage[active]{srcltx}
\usepackage{graphicx}
\usepackage{pdflscape}
\usepackage{amsmath}
\usepackage{amsthm}
\usepackage{lscape}
\usepackage{enumerate}
\usepackage{float}
\usepackage{grffile}
\usepackage{relsize}
\usepackage{tablefootnote}
\usepackage{footnote}
\usepackage{sectsty}
\usepackage{xcolor}
\usepackage{caption}
\usepackage{color}
\usepackage{natbib}
\usepackage{multirow}
\usepackage{bbm}
\usepackage[T1]{fontenc}
\usepackage[multiple]{footmisc}
\usepackage{appendix}
\usepackage{times}
\usepackage{relsize}
\usepackage{caption}
\usepackage{newtxtext,newtxmath}
\usepackage{enumitem}


%%%%%%%%%%%%%%%%%%%%%%%%%%%%%%%%%%%%%%%%%%%%%%%%%%%%%%%%%%%%
%%%   Title Page Information
%%%%%%%%%%%%%%%%%%%%%%%%%%%%%%%%%%%%%%%%%%%%%%%%%%%%%%%%%%%%
\title{ECON311: Assignment I}
\author{Leandro Freylejer}
\date{Due date: February 13th, 2026 before 4pm}

%%%%%%%%%%%%%%%%%%%%%%%%%%%%%%%%%%%%%%%%%%%%%%%%%%%%%%%%%%%%
%%%   Double and Single Space Commands
%%%%%%%%%%%%%%%%%%%%%%%%%%%%%%%%%%%%%%%%%%%%%%%%%%%%%%%%%%%%
\newlength{\defbaselineskip}
\setlength{\defbaselineskip}{\baselineskip}
\newcommand{\setlinespacing}[1]{\setlength{\baselineskip}{#1 \defbaselineskip}}
\newcommand{\doublespacing}{\setlength{\baselineskip}{1.3 \defbaselineskip}}
\newcommand{\singlespacing}{\setlength{\baselineskip}{1\defbaselineskip}}
\renewcommand{\thesubsection}{\arabic{subsection}}
\newtheoremstyle{dotlessS}{}{}{\itshape}{}{\color{dg}\bfseries}{}{ }{}
\theoremstyle{dotlessS}
\newtheorem{theorem}{Theorem}
\newtheorem{lemma}{Lemma}
\newtheorem{proposition}{Proposition}
\newtheorem{assumption}[theorem]{Assumption}
\newtheorem{corollary}{Corollary}
\newtheorem{definition}{Definition}
\newenvironment{example}[1][Example]{\begin{trivlist}
\item[\hskip \labelsep {\bfseries #1}]}{\end{trivlist}}
\newenvironment{remark}[1][Remark]{\begin{trivlist}
\item[\hskip \labelsep {\bfseries #1}]}{\end{trivlist}}

\newcommand\Pitem{%
  \addtocounter{enumi}{-1}%
  \renewcommand\theenumi{\arabic{\enumi}'}%
  \item%
  \renewcommand\theenumi{\arabic{enumi}}%
}

%%%%%%%%%%%%%%%%%%%%%%%%%%%%%%%%%%%%%%%%%%%%%%%%%%%%%%%%%%%%
%%%   Set Margins Here
%%%%%%%%%%%%%%%%%%%%%%%%%%%%%%%%%%%%%%%%%%%%%%%%%%%%%%%%%
\textwidth=6.5in
\textheight=9in
\oddsidemargin=0.10in
\evensidemargin=0.10in
\topmargin=-0.5in
%\parskip=3pt

%%%%%%%%%%%%%%%%%%%%%%%%%%%%%%%%%%%%%%%%%%%%%%%%%%%%%%%%%%%%%
%%%   Actual Document Begins Here
%%%%%%%%%%%%%%%%%%%%%%%%%%%%%%%%%%%%%%%%%%%%%%%%%%%%%%%%%%%%%
\begin{document}
\pagenumbering{arabic}
\maketitle
\vspace{-.3in}
\textbf{Instructions:} There are $5$ questions for a total of $80$ points. Answer all questions. To receive full marks, write your answers as complete and clear as possible. 

\section*{Question I [10 points]}
Complete the following table:
\begin{table}[H]
\begin{minipage}[h]{1\textwidth}
\centering
\resizebox{18cm}{!}{%
\begin{tabular}{|l|l|l|l|l|l|l|l|l|}

\hline
			&     &   	     &            &        &          &                                     &                              &                   \\[1.2ex]
Year            &Price of     &Quantity of     &Price of               &Quantity of       &RGDP Chained Weighted           &$\%\Delta RGDP_{F}$ & $GDP^{deflator}$ & Inflation \\
			&Burgers    &Burgers  	     &Magazines            &Magazines        &Prices  benchmark to 2014         &                                     &                              &  Rate                  \\[2.4ex]
\hline
			&     &   	     &            &        &          &                                     &                              &                   \\[1.2ex]
2013	       &$\$4.00$ &100                 &$\$2.00$               &180				&                                                 &                                     &                              &                   \\[2.4ex]
2014	       &$\$5.00$	&120		     &$\$2.50$	          &200      		      &                                                 &                                     &                              &                   \\[2.4ex]
2015             &$\$6.00$	&150		     &$\$3.50$                &200		       &                                                 &                                     &                              &                   \\[2.4ex]
\hline
\end{tabular}}
\end{minipage} 
\end{table}

\section*{Question II}
For this question you will need to download four data series. You can find all data at the World Bank Development Indicators website: \\
\textcolor{blue}{https://databank.worldbank.org/source/world-development-indicators}. The series needed to answer this question are: i. ``GDP per capita (constant 2015 $US\$$)'', ii. ``GDP per capita, PPP (constant 2021 international $\$$)''. You need to obtain these data series for both Canada and Chile from 1993 to 2017. To analyze the data you are free to use any software you would like. I recommend that you use Microsoft Excel because you can directly download the series in Excel format. If you do not have Excel in your personal computer, you can access it through the University's server.
\begin{enumerate}[label=\alph*]
\item (10 points) Plot GDP per capita (constant 2015 $US\$$) series for both countries on the same graph with time on the horizontal axis and GDP per capita on the vertical axis. On a different graph, plot GDP per capita, PPP (constant 2021 international $\$$) series for both countries with time on the horizontal axis and GDP per capita on the vertical axis. Which series is better at capturing differences in well-being across countries? Explain.
\item (10 points) Assume that in any given year, the GDP of a country is given by the production function studied in class ($Y=\overline{A}(K)^{\frac{1}{3}}(L)^{\frac{2}{3}}$). Also, as in class, assume that there are no differences in $\overline{A}$ across countries. Using the GDP per capita, PPP, calculate the annual differences in capital per capita for each year. Plot the amount of capital per capita in Canada relative to the amount of capital per capita in Chile for each year. Given our discussion in class, does this measure represent the actual differences in capital per capita? Explain.      
\item (10 points) Using  GDP per capita, PPP (constant 2021 international $\$$) calculate the average annual growth rates for the Chilean and Canadian economies for the entire period. Assume that the Chilean economy keeps growing at a constant rate equal to the average annual rate for the 1993-2017 period. How many years from 2017 would it take for the GDP per capita in Chile to reach the level of GDP per capita in Canada in 2017? 
\end{enumerate}  
\section*{Question III}
Suppose that total GDP in a country at a given year is given by the following production function:
\begin{eqnarray*}
Y=\overline{A}(K)^{\frac{1}{3}}(\beta L)^{\frac{2}{3}} 
\end{eqnarray*}
where $\beta$ represents labour productivity and all other variables are the same as in class. The introduction of $\beta$ expands the production model by allowing production to also depend on the human capital or productivity of the labour force.  Assume that there are two countries with the same number of workers ($\overline{L}$) and total factor productivity ($\overline{A}$) but with different levels of labour productivity ($\beta$). In particular, we are going to assume that the labour force in country 1 is twice as productive as the the labour force in country 2.
\begin{enumerate}[label=\alph*]
\item (5 points) Calculate the relative difference in GDP in country 1 relative to country 2 $\left(\frac{Y_1}{Y_2}\right)$ as a function of relative differences in the amount of physical capital  across countries $\left(\frac{K_1}{K_2}\right)$. In a graph with relative GDP differences $\left(\frac{Y_1}{Y_2}\right)$ on the vertical axis and relative differences in physical capital $\left(\frac{K_1}{K_2}\right)$ on the horizontal axis, plot this function. In the same graph, show what happens when there is an increase in the labour productivity of country 1 relative to the labour productivity of country 2.
\item (5 points) Assume that, in both countries, $\overline{L}$ grows at a constant rate $g_L$, $\overline{K}$ grows at a constant rate $g_K$, $\overline{A}$ grows at a constant rate $g_A$ and $\beta$ grows at a constant rate $g_{\beta}$. Calculate the constant rate of growth of GDP, $g_Y$, and GDP per capita, $g_y$.     
\end{enumerate}

\section*{Question IV [15 points]}
Consider a one time change in immigration policy that immediately and permanently increases the level of the labour force in an economy. Suppose it rises permanently from $\overline{L}$ to $\overline{L}'$. Assuming the economy starts in its initial steady state, use the Solow model to explain what happens to the economy over time (transition dynamics) and in the long run (steady state). In particular, discuss what happens to the levels of capital, capital per worker, GDP and GDP per worker both at the new steady state as well as during the transition to the new steady state. You are \textbf{not} required to discuss how growth rates change over time (this is tricky). To answer this question use both the Solow diagram with total output, and diagrams with time on the horizontal axis and each of these variables on the vertical axis.  

\section*{Question V}
Consider a Solow model where the production function no longer exhibits diminishing returns to capital accumulation. Assume that the production function is now $Y_t=\overline{A}K_t$. The rest of the model is unchanged.
\begin{enumerate}[label=\alph*]
\item (5 points) Draw the Solow diagram in this case.
\item (5 points) Suppose the economy begins with capital $\overline{K}_0$, show how the economy evolves over time.
\item (5 points) What happens to the growth rate of per capita GDP over time?
\end{enumerate}




\end{document}















\end{document}